% =====================================================================
\section{Dataset Preparation and Management}
\label{sec:dataset}
% =====================================================================

\subsection{The GRID audiovisual sentence corpus}
The proposal selects the GRID corpus as the primary dataset \spec. GRID
contains audio and frontal-face video recordings of 34 talkers, each speaking
1000 short sentences~\cite{cooke2006}; the proposal quotes approximately
33{,}000 usable sentences. Every sentence follows the same six-word grammar,
\begin{center}
\small\texttt{command}\;\;\texttt{colour}\;\;\texttt{preposition}\;\;\texttt{letter}\;\;\texttt{digit}\;\;\texttt{adverb},
\end{center}
as in ``place green at B 4 now''~\cite{cooke2006}. The vocabulary is therefore
small and closed, which makes GRID suitable for a six-week project but also
limits how far conclusions generalise to open conversation (\cref{sec:risks}).
Each utterance is distributed with a word-level alignment file giving start
and end times, from which the transcript is derived.

\subsection{Acquisition and organisation}
The team will download the video, audio and alignment archives for every
speaker from the official distribution and record the download date and
archive checksums in \texttt{data/raw/SOURCES.md}. The raw files are never
modified. All derived data are written to a separate tree:

\begin{lstlisting}[language={},caption={Proposed data directory layout \rec.},label={lst:layout}]
data/
  raw/grid/s1/video/*.mpg  s1/audio/*.wav  s1/align/*.align
  raw/musan/noise/...
  processed/
    mouths/s1/<utt_id>.npy        # uint8 [T, 96, 96]
    audio16k/s1/<utt_id>.wav      # 16 kHz mono, clean
    noisy_test/<snr>/<utt_id>.wav # fixed noisy copies for evaluation
  manifest.csv                     # one row per utterance
  splits.csv                       # speaker -> train / val / test
  stats/feature_norm.npz           # training-set mean and std
\end{lstlisting}

The manifest is the single index of the dataset. Each row records the
utterance identifier, speaker, transcript, number of video frames, audio
duration, the fraction of frames with a successful face detection, and the
paths of all derived files. Scripts select data by filtering the manifest,
never by listing directories.

\subsection{Transcription normalisation and vocabulary}
Transcripts are built from the alignment files by removing the silence and
short-pause tokens and joining the remaining words with single spaces. Text is
lower-cased and restricted to the characters \texttt{a}--\texttt{z} and space
\rec. The letter slot is spelled as a single character (``b'') and digits as
words (``four''), matching the alignment files. The output vocabulary is
therefore
\begin{equation}
\mathcal{V} = \{\texttt{<blank>}, \texttt{\textvisiblespace}, \texttt{a}, \dots, \texttt{z}\},
\qquad |\mathcal{V}| = 28,
\label{eq:vocab}
\end{equation}
with the CTC blank at index~0. A character vocabulary is chosen over a word
vocabulary because it keeps CER meaningful, lets the same model architecture
be reused for other text later, and matches the CTC systems cited in the
proposal~\cite{kanavnet2026,afouras2022}.

\subsection{Audio and video standardisation}
\begin{itemize}
  \item \textbf{Audio.} All audio is converted to 16\,kHz, 16-bit, mono WAV
        \rec. 16\,kHz is the standard rate for speech recognition, preserves the
        frequency range relevant to speech, and matches the rate the
        microphone path will deliver in the demonstration. Resampling uses a
        band-limited (anti-aliasing) resampler.
  \item \textbf{Video.} GRID video is recorded at 25 frames per second, and a
        3\,s utterance nominally yields 75 frames. Frames are decoded with
        OpenCV and their count checked against the audio duration. The training
        pipeline assumes exactly 25\,fps; the live pipeline resamples camera
        frames to 25\,fps by timestamp (\cref{sec:application}).
  \item \textbf{Synchronisation.} Audio and video of a GRID utterance start
        together. Sample $n$ of the 16\,kHz waveform corresponds to video frame
        $\lfloor 25n/16000 \rfloor$, so one video frame spans 640 audio samples
        and four 10\,ms audio feature frames. This fixed 4:1 ratio is used
        throughout \cref{sec:audio,sec:model}.
\end{itemize}

\subsection{Training, validation and test partitions}
The proposal requires evaluation on held-out speakers \spec. Partitions are
therefore made by speaker, not by utterance: every utterance from a given
speaker belongs to exactly one partition. The recommended split \rec\ is
\begin{itemize}
  \item \textbf{Test:} 4 speakers (2 female, 2 male), used only for final reporting;
  \item \textbf{Validation:} 3 speakers, used for early stopping and model selection;
  \item \textbf{Training:} all remaining speakers (about 27).
\end{itemize}
Speakers with missing or badly corrupted video, as identified by the integrity
checks below, are excluded before the split is drawn. The split is generated
once with a fixed seed, reviewed for gender balance, written to
\texttt{splits.csv} and committed to the repository.

\paragraph{Why speaker leakage matters.}
If utterances from one speaker appear in both training and test data, the
model can partly memorise that speaker's lip shape, face and voice. Test error
then measures recognition of familiar people rather than generalisation to new
users---the situation that matters for a captioning tool. Speaker leakage
would make every reported WER optimistically biased. Three safeguards prevent
it: the split file is the only source of partition membership; the data loader
raises an error if a speaker appears in more than one partition; and an
automated test (\cref{sec:testing}) asserts that the three speaker sets are
disjoint. Normalisation statistics are computed on the training partition only,
for the same reason.

\subsection{Dataset integrity checks}
A script \texttt{verify\_dataset.py} runs before any training and writes a
report. It checks that:
\begin{enumerate}
  \item every utterance has a video file, an audio file and an alignment file;
  \item video decodes without error and has $75 \pm 2$ frames;
  \item audio duration matches video duration to within one frame (40\,ms);
  \item the transcript contains exactly six words drawn from the GRID word list;
  \item Face Mesh succeeds on at least 90\% of frames \rec; utterances below
        this threshold are flagged and excluded;
  \item no audio file is silent or clipped throughout (RMS and peak checks).
\end{enumerate}
The report records how many utterances pass each check. These counts become
the dataset statistics quoted in the final report; no statistics are assumed
in advance.

\subsection{Noise preparation with MUSAN}
\label{sec:noise}
The proposal adds synthetic noise from the MUSAN corpus~\cite{snyder2015} at
five conditions: clean, 10\,dB, 5\,dB, 0\,dB and $-5$\,dB SNR \spec. MUSAN
contains separate music, speech and noise subsets. The core experiments use the
\emph{noise} subset \rec, which matches the proposal's ``background noise''
scenario; babble noise built from the speech subset is a possible later
experiment.

\paragraph{Noise partitions.} MUSAN noise files are split into disjoint
training, validation and test pools by file \rec. Using different noise
recordings at test time ensures that the evaluation measures robustness to
unseen noise rather than recognition of memorised noise patterns.

\paragraph{Mixing procedure.}
Let $s[n]$ be the clean speech waveform of length $N$ and $\tilde{w}[n]$ a noise
segment of the same length, cut at a random offset from a noise file (looped if
the file is shorter than $N$). The signal powers are
\begin{equation}
P_s = \frac{1}{N}\sum_{n=1}^{N} s[n]^2, \qquad
P_w = \frac{1}{N}\sum_{n=1}^{N} \tilde{w}[n]^2 .
\label{eq:power}
\end{equation}
To obtain a target SNR of $\rho$\,dB, the noise is scaled by
\begin{equation}
\alpha = \sqrt{\frac{P_s}{P_w \cdot 10^{\rho/10}}},
\qquad
y[n] = s[n] + \alpha\,\tilde{w}[n],
\label{eq:mix}
\end{equation}
which gives $10\log_{10}\!\big(P_s / (\alpha^2 P_w)\big) = \rho$ exactly. If
the peak $\max_n |y[n]|$ exceeds 1.0, the whole mixture $y$ is scaled down;
because speech and noise are scaled together, the SNR is unchanged.
\Cref{alg:mix} summarises the procedure.

\begin{algorithm}[H]
\caption{Mixing noise into clean speech at a target SNR.}
\label{alg:mix}
\begin{algorithmic}[1]
\Require clean speech $s$ (length $N$), noise pool $\mathcal{W}$, target SNR $\rho$ (dB), random generator $r$
\Ensure noisy speech $y$, measured SNR $\hat{\rho}$
\If{$\rho = \text{clean}$} \Return $s$, $\infty$ \EndIf
\State $w \gets$ file drawn from $\mathcal{W}$ using $r$
\State $\tilde{w} \gets$ segment of length $N$ from $w$ at random offset (loop $w$ if shorter)
\State $P_s \gets \operatorname{mean}(s^2)$;\quad $P_w \gets \operatorname{mean}(\tilde{w}^2)$
\State $\alpha \gets \sqrt{P_s / (P_w \cdot 10^{\rho/10})}$
\State $y \gets s + \alpha\tilde{w}$
\State $\hat{\rho} \gets 10\log_{10}\!\big(P_s / \operatorname{mean}((\alpha\tilde{w})^2)\big)$ \Comment{verification}
\State \textbf{assert} $|\hat{\rho} - \rho| < 0.01$
\If{$\max|y| > 1$} $y \gets y / \max|y|$ \Comment{preserves SNR}\EndIf
\State \Return $y$, $\hat{\rho}$
\end{algorithmic}
\end{algorithm}

\paragraph{Speech power definition.}
\Cref{eq:power} measures speech power over the whole clip, including the short
silences at the start and end of each GRID utterance \rec. This definition is
simple and reproducible. Its drawback is that the SNR during active speech is
somewhat higher than the nominal value. The final report will state this
definition explicitly; measuring power over voiced frames only is an
alternative that the team may adopt if consistency with a particular reference
paper becomes important \tbd.

\paragraph{Training versus test noise.}
During training, noise is mixed on the fly: for each utterance an SNR is drawn
from $\{\text{clean}, 10, 5, 0, -5\}$\,dB with equal probability \rec, so the
model sees each utterance under many conditions over the course of training.
For validation and testing, each utterance is mixed once per SNR with a fixed
seed and the result is stored in \texttt{processed/noisy\_test/}. Every model
is then evaluated on exactly the same audio.

\paragraph{Labels under noise.}
Adding noise changes only the audio waveform. The video and transcript of the
utterance are unchanged, so every noisy copy keeps the transcript of its clean
original. This property is what allows the same test sentences to be scored at
every SNR and the results compared directly.

\subsection{Storage of processed features and metadata}
Mouth crops are stored as one uncompressed \texttt{uint8} NumPy array per
utterance \rec, which allows fast memory-mapped loading; grey-scale storage
reduces size threefold compared with colour. Clean audio is stored as 16\,kHz
WAV, and log-Mel features are computed on the fly on the GPU, because the
noise---and therefore the features---changes every time an utterance is
loaded during training. Feature normalisation statistics (per-band mean and
standard deviation of the log-Mel features over clean and noisy training audio)
are stored in \texttt{stats/feature\_norm.npz} and loaded by both training and
inference code.
